\section*{Введение}
Работа посвящена статье Christopher A. Sims и Tao Zha «Does Monetary Policy Generate Recessions?» \cite{SIMS2006}. Главный исследовательский вопрос статьи состоит в том, вызывают ли денежно-кредитные шоки рецессии или наблюдаемые спады в основном отражают систематическую реакцию денежно-кредитной политики на состояние экономики.

Авторы рассматривают рецессии в США и пытаются отделить неожиданные шоки денежно-кредитной политики от реакции ФРС на текущую макроэкономическую ситуацию. Для этого они строят структурную векторную авторегрессию (SVAR) по квартальным данным и используют ограничения, которые позволяют выделить денежно-кредитные шоки из совместной динамики выпуска, цен, денежной массы, процентных ставок и других переменных.

Теоретическая мотивация модели опирается на динамическую стохастическую модель общего равновесия (DSGE). Авторы задают денежно-кредитный шок внутри DSGE-модели, получают данные из этой модели, а затем проверяют, можно ли восстановить этот шок по ограниченному набору наблюдаемых переменных. Для этого используется фильтр Калмана: если остаточная дисперсия шока после проекции на наблюдаемые переменные мала, то шок можно считать достаточно хорошо восстановимым. На этом основании авторы обосновывают, что их SVAR-идентификация может выделять денежно-кредитные шоки из эмпирических данных.

Отдельное внимание в статье уделяется критике части литературы. Авторы считают слишком сильным предположение, что большие блоки переменных частного сектора вообще не реагируют мгновенно на денежно-кредитную политику. Например, в некоторых работах долгосрочная ставка или товарные цены фактически рассматривались как переменные, которые не могут измениться в тот же период после монетарного шока. Sims и Zha считают такие ограничения экономически сомнительными, поэтому разрешают одновременные связи между финансовыми переменными и частью товарных цен.

Также авторы оценивают несколько спецификаций модели. Это связано с тем, что в литературе не было единого мнения о том, какой показатель лучше отражает денежно-кредитную политику: денежный агрегат M2 или показатель, связанный с резервами банков (Total Reserves). Поэтому в статье рассматриваются обе версии модели и проверяется, насколько выводы о роли денежно-кредитных шоков устойчивы к выбору денежного индикатора. В этих версиях используются разные процентные ставки: в модели с M2 — 90-дневная ставка по казначейским векселям, а в модели с резервами — ставка федеральных фондов.

Основной вывод статьи состоит в том, что неожиданные денежно-кредитные шоки играли ограниченную роль в объяснении послевоенных рецессий в США. Большая часть изменений в денежных переменных отражает не случайные сдвиги политики, а систематическую реакцию ФРС на состояние экономики.
